% Abschnitt 2.2: Modell für die Ladezeit (mit Gleichung und Aufzählung)
\subsection{Modell für die Ladezeit}
\label{sec:modell}

Fasst man die Schritte aus Abbildung~\ref{fig:ablauf} zusammen, ergibt sich
für die Ladezeit $t_\mathrm{L}$ näherungsweise
\begin{equation}
  t_\mathrm{L} = t_\mathrm{DNS} + t_\mathrm{V} + t_\mathrm{S} + \frac{S}{B}\,.
  \label{eq:ladezeit}
\end{equation}
Dabei ist $t_\mathrm{DNS}$ die Dauer der Namensauflösung, $t_\mathrm{V}$ die
Dauer des Verbindungsaufbaus, $t_\mathrm{S}$ die Bearbeitungszeit des Servers,
$S$ die übertragene Datenmenge und $B$ die Datenrate der Verbindung. Das
Modell vernachlässigt parallele Übertragungen und die Zeit für das Rendern im
Browser, zeigt aber, wo sich ansetzen lässt:
\begin{itemize}
  \item Datenmenge $S$ verringern, etwa durch Komprimierung, moderne
        Bildformate und das Entfernen ungenutzter Skripte.
  \item Wiederholte Übertragungen vermeiden, indem der Browser unveränderte
        Dateien aus dem Cache lädt.
  \item Wartezeiten $t_\mathrm{V}$ und $t_\mathrm{S}$ senken, zum Beispiel
        durch Server in der Nähe der Nutzer und durch vorab erzeugte Seiten.
\end{itemize}
